\section{Estado del Arte}

En los últimos años, el campo de la inteligencia artificial aplicada a la predicción de complicaciones postoperatorias ha experimentado un crecimiento significativo, siguen siendo pocos los estudios que destacan por sus resultados o procedimientos en la prevención del shock hemorrágico. Los modelos de ML surgen entonces como herramientas prometedoras para superar las limitaciones de los métodos tradicionales al integrar múltiples variables clínicas y proporcionar predicciones más precisas. Sin embargo, aún quedan desafíos metodológicos y prácticos que limitan su implementación clínica.

\subsection{Modelo predictivo para hemorragia pospancreatectomía}

Zhang et al.\ \cite{zhang2025} desarrollaron y validaron externamente un modelo de machine learning para la estratificación precisa del riesgo de hemorragia pospancreatectomía (PPH) mediante una cohorte internacional. El estudio utilizó una cohorte de entrenamiento de 3,609 pacientes de la base de datos del American College of Surgeons National Surgical Quality Improvement Program (ACS-NSQIP) y una cohorte de validación externa de 1,347 pacientes del National Cancer Center en China. Los investigadores compararon múltiples algoritmos evaluados, incluyendo GBM, Deep Neural Networks (DNN) y Logistic Regression, seleccionando finalmente el Lasso-GBM, el cual identificó predictores claves como: el índice de masa corporal (IMC), la clasificación ASA, la administración de esteroides preoperatorios, el uso de betabloqueadores, el tiempo quirúrgico y los niveles de hematocrito. Este enfoque usa regresión LASSO (Least Absolute Shrinkage and Selection Operator). Esta técnica es un modelo de IA que iterativamente penaliza las variables no esenciales, resuelve problemas de multicolinealidad y previene el overfitting. Sin embargo, presentó limitaciones importantes por la falta de datos temporales como el momento en el que comenzó la hemorragia y la dependencia de protocolos de centros de alta volumetría limitó su aplicabilidad en entornos menos estandarizados.

\subsection{Modelo predictivo para accidente cerebrovascular perioperatorio}

Cong et al.\ \cite{cong2025} desarrollaron y validaron un modelo basado en machine learning para predecir el riesgo de accidente cerebrovascular (stroke) perioperatorio en pacientes sometidos a cirugías no cardíacas, no vasculares y no neuroquirúrgicas. Utilizaron datos de 2,986 pacientes del Hospital Provincial de Henan, comparando nueve algoritmos diferentes de machine learning. El estudio desembocó en el uso de los modelos de LightGBM y KNN mostraron el mejor desempeño, con AUCs de 0.941 y 0.988 respectivamente en el conjunto de entrenamiento. Este estudio destacó por su manejo del desbalance de clases (solo el 16-20\% de los pacientes experimentaron stroke), utilizando métricas clínicas relevantes como sensibilidad y F1 Score. Además, incorporó SHAP para explicar las predicciones del modelo, facilitando la interpretación clínica. Sin embargo, el estudio y su enfoque en predicción preoperatoria no abordó la detección temprana postoperatoria que sería necesaria para intervenciones inmediatas.

\subsection{Predicción del recuperamiento de shock hemorrágico severo usando regresión logística}

Aunque la mayoría de los avances recientes en predicción de shock hemorrágico se han centrado en modelos clínicos con machine learning, algunos estudios preclínicos han explorado enfoques similares en modelos animales. En este caso, Lucas et al.\ \cite{lucas2019} desarrollaron un modelo de regresión logística con selección rigurosa de características que logró predecir la recuperación tras shock hemorrágico severo en ratas con una precisión del 84\%, superando a los métodos tradicionales basados únicamente en presión arterial y frecuencia cardíaca. Este enfoque metodológico, aunque no directamente aplicable en entornos clínicos humanos, refuerza el potencial de los modelos predictivos basados en múltiples parámetros fisiológicos para mejorar la toma de decisiones en situaciones críticas, lo que respalda la viabilidad de desarrollar sistemas similares adaptados al entorno postoperatorio humano.

\subsection{Modelos multimodales para pronóstico en UCI}

Pal et al.\ \cite{pal2022} presentaron ShockModes, un modelo multimodal para la predicción de resultados en UCI que integra tres fuentes de datos heterogéneas: notas clínicas de médicos procesadas con DocBERT (un Transformer entrenado en textos médicos), series temporales de signos vitales (frecuencia cardíaca, presión arterial, SpO\textsubscript{2}) procesadas con Regresión Logística, Random Forest, GradientBoost, AdaBoost y XGBoost, y datos de comorbilidades del paciente. El modelo alcanzó un recall de 0.81 y una especificidad de 0.94, demostrando el potencial de las arquitecturas multimodales en este dominio. Sin embargo, los propios autores reconocen que este enfoque requiere volúmenes de datos significativamente mayores y una infraestructura de recolección de datos más compleja que la disponible en la mayoría de centros clínicos. En nuestro caso, contamos únicamente con una de estas fuentes de datos (variables tabulares), lo que contextualiza las limitaciones de nuestro modelo en comparación con este referente.

\subsection{Limitaciones de los modelos predictivos clínicos}

Ma et al.\ \cite{ma2023} realizaron una revisión sistemática sobre el manejo deficiente de predictores continuos en modelos de predicción clínica. Su análisis reveló que la dicotomización de variables continuas (como convertir la edad o la hemoglobina en categorías binarias) y el uso de categorías arbitrarias para predictores continuos son prácticas extendidas que reducen significativamente la precisión predictiva de los modelos, al disminuir la información disponible sobre la verdadera distribución de cada variable. Esto es relevante para nuestro trabajo: los experimentos de categorización de EDAD y HB\_PREQX no produjeron mejoras y, de acuerdo con estos hallazgos, pudieron haber introducido pérdida de información.

Complementariamente, Steyerberg et al.\ \cite{steyerberg2018} analizaron las principales problemáticas de los enfoques de modelado tradicionales en contextos clínicos, identificando como riesgos críticos: el uso de datasets de tamaño reducido esperando generalización amplia, la categorización de predictores continuos, y la sobreestimación de los resultados del análisis estadístico (modelos que presentan resultados muy favorables en la evaluación interna pero que rinden significativamente peor en aplicación clínica real). Este último punto es especialmente pertinente en nuestro estudio, donde la optimización del umbral de decisión permitió mejorar el rendimiento en el conjunto de prueba a pesar de que los valores de la validación cruzada fueron moderados.

\subsection{Síntesis y brecha de conocimiento}

La revisión de estos estudios revela importantes avances metodológicos en la aplicación de machine learning para predecir complicaciones postoperatorias, pero también identifica una brecha crítica que aborda este proyecto: la ausencia de un sistema especializado para la detección temprana del shock hemorrágico postoperatorio en cirugía mayor no cardíaca. Mientras Zhang et al.\ \cite{zhang2025} desarrollaron un modelo predictivo para hemorragia pospancreatectomía y Cong et al.\ \cite{cong2025} se enfocaron en la predicción de stroke perioperatorio, ninguno aborda específicamente la detección temprana del shock hemorrágico en el período postoperatorio, cuando los signos clínicos pueden ser sutiles pero la intervención oportuna es crítica para la supervivencia. El estudio ShockModes \cite{pal2022} demuestra el potencial de las arquitecturas multimodales, pero requiere fuentes de datos heterogéneas que superan los datos disponibles en este proyecto. La brecha que aborda este proyecto es, por tanto, muy específica: desarrollar un sistema de alerta temprana basado en inteligencia artificial que permita identificar pacientes en riesgo de shock hemorrágico durante el período postoperatorio, reconociendo las limitaciones metodológicas señaladas por Ma et al.\ \cite{ma2023} y Steyerberg et al.\ \cite{steyerberg2018}.

\begin{table}[H]
\centering
\caption{Estado del arte en modelos predictivos para complicaciones postoperatorias}
\label{tab:estado-arte}
\small
\begin{tabular}{|p{2.5cm}|p{3cm}|p{4.5cm}|p{4.5cm}|}
\hline
\textbf{Proyecto} & \textbf{Temática principal} & \textbf{Metodología} & \textbf{Limitaciones} \\ \hline
Zhang et al.\ \cite{zhang2025} & Modelo predictivo para hemorragia pospancreatectomía & Modelo Lasso-GBM validado externamente; enfoque en explicabilidad clínica & Diseño retrospectivo restringió acceso a variables quirúrgicas detalladas; dependencia de protocolos de centros de alta volumetría \\ \hline
Cong et al.\ \cite{cong2025} & Modelo predictivo para stroke perioperatorio & LightGBM y KNN mostraron mejor desempeño (AUC 0.941 y 0.988); utilizó SHAP para explicabilidad clínica & Diseño retrospectivo con acceso limitado a variables quirúrgicas granulares; enfoque en predicción preoperatoria, no en detección temprana postoperatoria; no específico para shock hemorrágico \\ \hline
Lucas et al.\ \cite{lucas2019} & Predicción de recuperación tras shock hemorrágico severo mediante regresión logística & Regresión logística con selección rigurosa de características; evaluación de tres clasificadores: con todas las características, con conjunto reducido y solo MAP y HR & Precisión del 84\% en predicción de recuperación; modelo aplicado en contexto preclínico (ratas); no generalizable directamente a entornos clínicos humanos \\ \hline
Pal et al.\ \cite{pal2022} & Modelo multimodal para pronóstico en UCI (ShockModes) & Integración de notas clínicas (DocBERT), series temporales de signos vitales y comorbilidades; recall 0.81, especificidad 0.94 & Requiere múltiples fuentes de datos heterogéneas; infraestructura de recolección compleja; necesita volúmenes de datos mayores según los autores \\ \hline
Ma et al.\ \cite{ma2023}; Steyerberg et al.\ \cite{steyerberg2018} & Limitaciones metodológicas en modelos predictivos clínicos & Revisión sistemática de prácticas de modelado en regresión logística y validación interna & Categorizar predictores continuos reduce la precisión predictiva; datasets pequeños no garantizan generalización; sobreestimación de resultados estadísticos respecto al rendimiento clínico real \\ \hline
\end{tabular}
\end{table}

\newpage

% -----------------------------------------------------------------------------
% 4. METODOLOGÍA
% -----------------------------------------------------------------------------
